\documentclass[11pt]{article}

\usepackage[utf8]{inputenc}
\usepackage[T1]{fontenc}
\usepackage[spanish,shorthands=off]{babel}
\usepackage[margin=1.9cm, top=1.6cm, bottom=1.8cm]{geometry}
\usepackage{amsmath,amssymb}
\usepackage{mathpazo}
\usepackage{xcolor}
\usepackage{tcolorbox}
\tcbuselibrary{skins}
\usepackage{tikz}
\usepackage{enumitem}
\usepackage{titlesec}
\usepackage{caption}
\usepackage{subcaption}
\usepackage{booktabs}
\usepackage{array}
\usetikzlibrary{patterns, arrows.meta, decorations.pathmorphing, calc}

\definecolor{ucred}{RGB}{117,26,41}
\definecolor{devcolor}{RGB}{45,70,95}
\definecolor{acero}{RGB}{120,130,145}
\definecolor{agua}{RGB}{40,110,180}

\titleformat{\section}{\normalfont\large\bfseries\color{ucred}}{\thesection}{0.5em}{}
\titlespacing*{\section}{0pt}{0.8em}{0.35em}
\setlength{\parskip}{2pt}
\setlist{leftmargin=1.2em, itemsep=1pt, topsep=2pt, parsep=0pt}
\setlength{\textfloatsep}{8pt}
\setlength{\floatsep}{6pt}
\setlength{\intextsep}{6pt}

\newtcolorbox{desarrollo}[1]{
  enhanced, colback=white, colframe=devcolor, coltitle=white, fonttitle=\bfseries,
  title={#1}, boxrule=0.8pt, arc=1.5mm, top=4pt, bottom=4pt, left=7pt, right=7pt
}

\captionsetup{font=small, labelfont={bf,color=ucred}, skip=3pt}
\tikzset{cota/.style={<->, >={Stealth[length=1.4mm]}, thin}}

\begin{document}

%================================================================
\begin{center}
  {\scshape\small Pontificia Universidad Cat\'olica de Chile \;$\cdot$\; ICM 3323 --- Elementos Finitos en Dise\~no Mec\'anico}\\[3pt]
  {\Large\bfseries\color{ucred} Optimizaci\'on del blindaje de acero de una compuerta de fondo}\\[2pt]
  {\normalsize Propuesta de proyecto T4 \;$\cdot$\; Pedro Pablo Taboada Santiba\~nez \;$\cdot$\; \today}\\[-2pt]
  {\color{ucred}\rule{\linewidth}{0.6pt}}
\end{center}
\vspace{-6pt}

%================================================================
\section{Contexto y motivaci\'on}
%================================================================

El desag\"ue de fondo del Proyecto Hidroel\'ectrico Los Lagos tiene dos compuertas deslizantes en
serie dentro de un conducto rectangular de $2300\times4600$~mm. En esa zona el conducto se reviste
con un \textbf{blindaje de acero}: planchas soldadas en forma de caj\'on, en tres tramos (1500, 1420 y
1500~mm, figura~\ref{fig:corte}), ancladas al hormig\'on con pernos M20. Los datos vienen del informe
de c\'alculo de Voith \emph{1530-VOI-320-0-CC-0072} y de los planos del proyecto.

\begin{figure}[h]
\centering
\begin{tikzpicture}[x=1.3cm, y=0.4cm, >={Stealth[length=1.8mm]}]
  \def\H{4.6}
  \fill[pattern=north east lines, pattern color=gray!70] (-0.2,-0.9) rectangle (8.6,0);
  \fill[pattern=north east lines, pattern color=gray!70] (-0.2,\H) rectangle (2.55,\H+0.9);
  \fill[pattern=north east lines, pattern color=gray!70] (2.95,\H) rectangle (4.37,\H+0.9);
  \fill[pattern=north east lines, pattern color=gray!70] (4.77,\H) rectangle (8.6,\H+0.9);
  \fill[agua!8] (-0.2,0) rectangle (8.6,\H);
  \foreach \xa/\xb in {1.05/2.55, 2.95/4.37, 4.77/6.27}{
    \draw[line width=2.2pt, devcolor] (\xa,0) -- (\xb,0) (\xa,\H) -- (\xb,\H);
  }
  \foreach \xg/\txt in {2.55/compuerta 1, 4.37/compuerta 2}{
    \fill[gray!25] (\xg,\H) rectangle (\xg+0.4,\H+1.8);
    \draw (\xg,\H) -- (\xg,\H+1.8) (\xg+0.4,\H) -- (\xg+0.4,\H+1.8);
    \fill[acero!70] (\xg+0.12,0) rectangle (\xg+0.28,\H);
    \draw[thick] (\xg+0.12,0) rectangle (\xg+0.28,\H);
    \node[font=\small, anchor=south] at (\xg+0.2,\H+1.8) {\txt};
  }
  \draw[thick] (-0.2,0) -- (8.6,0) (-0.2,\H) -- (2.55,\H) (2.95,\H) -- (4.37,\H) (4.77,\H) -- (8.6,\H);
  \draw[->, very thick, agua] (-0.05,\H/2) -- (0.9,\H/2) node[midway, above, font=\small] {flujo};
  \foreach \xa/\xb/\txt in {1.05/2.55/1500, 2.95/4.37/1420, 4.77/6.27/1500}{
    \draw[cota] (\xa,-1.9) -- (\xb,-1.9) node[midway, fill=white, inner sep=1pt, font=\small] {\txt};
    \draw[thin, gray] (\xa,-0.1) -- (\xa,-2.2) (\xb,-0.1) -- (\xb,-2.2);
    \node[font=\footnotesize, text=devcolor] at ({(\xa+\xb)/2},\H/2) {blindaje};
  }
  \draw[cota] (7.5,0) -- (7.5,\H) node[midway, fill=agua!8, inner sep=1pt, font=\small] {4600};
  \node[font=\scriptsize, anchor=west] at (8.6,0) {EL 78.01};
  \node[font=\scriptsize, anchor=west] at (8.6,\H) {EL 82.61};
  \node[font=\scriptsize, anchor=west, agua] at (8.6,\H+2.2) {NA 111.0};
\end{tikzpicture}
\caption{Corte longitudinal esquem\'atico del desag\"ue de fondo (medidas en mm, no a escala).}
\label{fig:corte}
\end{figure}

El caso que dimensiona el blindaje es el \textbf{conducto vac\'io}: el agua filtrada queda detr\'as
de las planchas con la presi\'on del embalse, $p(y) = \gamma_w\,(111.0 - \text{cota})$, de 2.84 a
3.30~kgf/cm$^2$, y las empuja hacia el conducto. El hormig\'on no puede tirar de la plancha, as\'i que
\'esta flecta como placa apoyada solo en los pernos. El dise\~no original (plancha de 20~mm, pernos a
$\le300$~mm, unos 290 pernos por tramo) da una tensi\'on m\'axima de unos 824~kgf/cm$^2$ frente a un
admisible de 2100~kgf/cm$^2$ (DIN~19704): la plancha trabaja \textbf{a menos del 40\,\%}. Hay margen
para ahorrar, pero espesor y pernos compiten: una plancha m\'as delgada pide pernos m\'as juntos, y
cada perno suma soldadura y montaje.

\begin{desarrollo}{Pregunta del proyecto}
?`Qu\'e combinaci\'on de \textbf{espesor de plancha} y \textbf{distribuci\'on de pernos} da el blindaje
con menos acero y menos pernos que cumpla las tensiones admisibles, la capacidad de los pernos y una
flecha razonable con el conducto vac\'io?
\end{desarrollo}

%================================================================
\section{Modelo y optimizaci\'on}
%================================================================

Un tramo se modela como un caj\'on de acero A572 Gr~50 ($E = 2.1\times10^6$~kgf/cm$^2$, $\nu=0.3$) de
ancho $a = 2300$~mm, alto $b = 4600$~mm, largo $L$ y espesor $t$ (figura~\ref{fig:modelo}a). La
presi\'on hidrost\'atica act\'ua en la cara exterior; cada perno es un apoyo de r\'otula
($\mathbf u = 0$, giro libre) y los extremos $z = 0, L$ se fijan, como en el informe original.
Las variables de dise\~no son el espesor $t$ y la distribuci\'on de pernos de cada plancha
(separaciones $p_s$, $p_z$ y distancia al borde $e$, figura~\ref{fig:modelo}b). En lugar de una
optimizaci\'on formal se hace un barrido acotado:
\begin{itemize}
  \item \textbf{Cuatro distribuciones de pernos} (figura~\ref{fig:distrib}): A, la grilla original a
        300~mm; B, una grilla m\'as abierta; C, tresbolillo; y D, grilla densificada cerca de las
        esquinas y los extremos, donde se concentra la carga.
  \item \textbf{Varios espesores comerciales}, por ejemplo $t = 12$, 14, 16, 18 y 20~mm (m\'as 2~mm de
        corrosi\'on).
\end{itemize}
Cada combinaci\'on se revisa contra: $\sigma_{VM} \le 2100$~kgf/cm$^2$ en la plancha; en el perno
m\'as cargado, tracci\'on $R/A_s$, soldadura $R/A_w$ y arranque del hormig\'on $R/(\pi D h)$ bajo sus
admisibles; flecha $w_{max} \le p_s/500$; y $e \le 125$~mm. De ah\'i se elige la \textbf{mejor
distribuci\'on} (la que reparte mejor la carga con menos pernos) y el \textbf{menor espesor} que cumple
con ella, y esa combinaci\'on se deja como el \'optimo. Si queda tiempo, se refina alrededor de ese
punto con una optimizaci\'on continua de $t$, $p_s$, $p_z$ y $e$.

\begin{figure}[h]
\centering
\begin{subfigure}[b]{0.42\linewidth}\centering
\begin{tikzpicture}[scale=0.75, x={(1cm,0cm)}, y={(0cm,1cm)}, z={(0.78cm,0.5cm)}, >={Stealth[length=1.6mm]}]
  \def\A{1.1}  \def\B{2.2}  \def\T{0.14}  \def\L{3.0}
  \pgfmathsetmacro{\Ao}{\A+\T} \pgfmathsetmacro{\Bo}{\B+\T}
  \fill[gray!8] (-\A,-\B,0) -- (-\A,-\B,\L) -- (-\A,\B,\L) -- (-\A,\B,0) -- cycle;
  \fill[gray!14] (-\A,-\B,0) -- (\A,-\B,0) -- (\A,-\B,\L) -- (-\A,-\B,\L) -- cycle;
  \draw[gray!60] (-\A,-\B,\L) -- (\A,-\B,\L) (-\A,-\B,\L) -- (-\A,\B,\L);
  \fill[acero!35] (\Ao,-\Bo,0) -- (\Ao,-\Bo,\L) -- (\Ao,\Bo,\L) -- (\Ao,\Bo,0) -- cycle;
  \fill[acero!20] (-\Ao,\Bo,0) -- (\Ao,\Bo,0) -- (\Ao,\Bo,\L) -- (-\Ao,\Bo,\L) -- cycle;
  \draw (\Ao,-\Bo,0) -- (\Ao,-\Bo,\L) -- (\Ao,\Bo,\L) -- (-\Ao,\Bo,\L) -- (-\Ao,\Bo,0);
  \draw (\Ao,\Bo,0) -- (\Ao,\Bo,\L);
  \foreach \zz in {0.35, 1.1, 1.9, 2.65}{
    \foreach \yy in {-1.95,-1.4,...,2.0}{ \fill[ucred] (\Ao,\yy,\zz) circle (1.4pt); }
    \foreach \xx in {-0.85,-0.3,0.25,0.8}{ \fill[ucred] (\xx,\Bo,\zz) circle (1.4pt); }
  }
  \fill[acero!70, even odd rule] (-\Ao,-\Bo,0) -- (\Ao,-\Bo,0) -- (\Ao,\Bo,0) -- (-\Ao,\Bo,0) -- cycle
                                  (-\A,-\B,0) -- (\A,-\B,0) -- (\A,\B,0) -- (-\A,\B,0) -- cycle;
  \draw (-\Ao,-\Bo,0) -- (\Ao,-\Bo,0) -- (\Ao,\Bo,0) -- (-\Ao,\Bo,0) -- cycle
        (-\A,-\B,0) -- (\A,-\B,0) -- (\A,\B,0) -- (-\A,\B,0) -- cycle;
  \foreach \yy in {-2.0,-1.2,-0.4,0.4,1.2,2.0}{
    \pgfmathsetmacro{\lp}{0.75 - 0.12*\yy}
    \draw[->, agua, thick] (\Ao+\lp,\yy,1.5) -- (\Ao+0.05,\yy,1.5);
  }
  \draw[agua] (\Ao+0.99,-2.0,1.5) -- (\Ao+0.51,2.0,1.5);
  \foreach \xx in {-0.6,0.5}{ \draw[->, agua, thick] (\xx,\Bo+0.5,2.0) -- (\xx,\Bo+0.05,2.0); }
  \node[agua, font=\small, anchor=west] at (\Ao+1.0,0.0,1.5) {$p(y)$};
  \draw[cota] (-\A,-\B-0.45,0) -- (\A,-\B-0.45,0) node[midway, below, font=\small] {$a$};
  \draw[cota] (-\A-0.5,-\B,0) -- (-\A-0.5,\B,0) node[midway, left, font=\small] {$b$};
  \draw[cota] (\Ao+0.25,-\Bo,0) -- (\Ao+0.25,-\Bo,\L) node[midway, below right, font=\small] {$L$};
  \draw[<-, thin] (-\A-0.07,1.2,0) -- ++(-0.8,0.5,0) node[left, font=\small] {$t$};
  \begin{scope}[shift={(-2.6,-2.6,0)}]
    \draw[->] (0,0,0) -- (0.7,0,0) node[right, font=\scriptsize] {$x$};
    \draw[->] (0,0,0) -- (0,0.7,0) node[above, font=\scriptsize] {$y$};
    \draw[->] (0,0,0) -- (0,0,0.9) node[above right, font=\scriptsize] {$z$};
  \end{scope}
\end{tikzpicture}
\caption{Tramo de blindaje con pernos (\textcolor{ucred}{$\bullet$}) y presi\'on de filtraci\'on.}
\end{subfigure}\hfill
\begin{subfigure}[b]{0.56\linewidth}\centering
\begin{tikzpicture}[x=0.7cm, y=0.7cm, >={Stealth[length=1.4mm]}]
  \fill[acero!15] (0,0) rectangle (9.2,4.5);
  \draw[thick] (9.2,0) -- (0,0) -- (0,4.5) -- (9.2,4.5);
  \draw[decorate, decoration={zigzag, segment length=5pt, amplitude=1.5pt}] (9.2,0) -- (9.2,4.5);
  \foreach \i in {0,...,6}{ \foreach \j in {0,...,3}{
      \fill[ucred] ({0.6+1.35*\i},{0.6+1.1*\j}) circle (1.8pt); } }
  \draw[cota] (0,-0.4) -- (0.6,-0.4) node[midway, below, font=\small] {$e$};
  \draw[cota] (0.6,-0.4) -- (1.95,-0.4) node[midway, below, font=\small] {$p_s$};
  \draw[cota] (1.95,-0.4) -- (3.3,-0.4) node[midway, below, font=\small] {$p_s$};
  \draw[thin, gray] (0.6,0.5) -- (0.6,-0.55) (1.95,0.5) -- (1.95,-0.55) (3.3,0.5) -- (3.3,-0.55);
  \draw[cota] (-0.4,0) -- (-0.4,0.6) node[midway, left, font=\small] {$e$};
  \draw[cota] (-0.4,0.6) -- (-0.4,1.7) node[midway, left, font=\small] {$p_z$};
  \draw[cota] (-0.4,1.7) -- (-0.4,2.8) node[midway, left, font=\small] {$p_z$};
  \draw[thin, gray] (0.5,0.6) -- (-0.55,0.6) (0.5,1.7) -- (-0.55,1.7) (0.5,2.8) -- (-0.55,2.8);
  \node[rotate=90, font=\small, anchor=south] at (-1.3,2.25) {$z$};
  \node[font=\small, anchor=west] at (3.6,-0.95) {$s$ (a lo largo de la pared)};
  % pano entre pernos, en corte
  \begin{scope}[shift={(0,-4.6)}]
    \fill[acero!60] (0,0) rectangle (9.2,0.3);
    \draw (0,0) -- (9.2,0) (0,0.3) -- (9.2,0.3);
    \foreach \xp in {0.6,4.6,8.6}{
      \draw[line width=1.6pt, ucred] (\xp,0.3) -- (\xp,1.3);
      \fill[ucred] (\xp-0.3,1.25) rectangle (\xp+0.3,1.4);
      \draw[ucred, fill=white, thick] (\xp,0.3) circle (1.8pt);
    }
    \foreach \xx in {1.2,1.8,...,4.2, 5.2,5.8,...,8.2}{ \draw[->, agua] (\xx,1.0) -- (\xx,0.33); }
    \draw[dashed, thick, devcolor] (0.6,0) .. controls (1.9,-1.2) and (3.3,-1.2) .. (4.6,0)
                                         .. controls (5.9,-1.2) and (7.3,-1.2) .. (8.6,0);
    \fill[ucred] (2.6,-0.9) circle (1.8pt) node[below, font=\small, black] {$\sigma_c$};
    \draw[->, thin] (4.6,-1.15) node[below, font=\small] {$\sigma_p$} -- (4.6,-0.1);
    \node[agua, font=\small, fill=acero!0, inner sep=1pt] at (6.7,1.25) {$p$};
  \end{scope}
\end{tikzpicture}
\caption{Arriba: grilla de pernos de una pared vista desde el hormig\'on. Abajo: corte del pa\~no
entre pernos, deformada y tensiones a controlar.}
\end{subfigure}
\caption{Modelo de un tramo y variables de dise\~no.}
\label{fig:modelo}
\end{figure}

\begin{figure}[h]
\centering
\begin{tikzpicture}[x=0.78cm, y=0.78cm]
  \tikzset{panel/.style={fill=acero!15, draw=black, thick},
           perno/.style={fill=ucred}}
  \def\W{4.0} \def\Hp{2.6}
  % A: grilla original
  \begin{scope}[shift={(0,0)}]
    \fill[acero!15] (0,0) rectangle (\W,\Hp); \draw[thick] (0,0) rectangle (\W,\Hp);
    \foreach \s in {0.3,0.9,...,3.8}{ \foreach \z in {0.3,0.8,...,2.4}{ \fill[perno] (\s,\z) circle (1.6pt); } }
    \node[font=\small, anchor=north, align=center] at (\W/2,-0.1) {\textbf{A}: original, 300~mm};
  \end{scope}
  % B: grilla abierta
  \begin{scope}[shift={(4.6,0)}]
    \fill[acero!15] (0,0) rectangle (\W,\Hp); \draw[thick] (0,0) rectangle (\W,\Hp);
    \foreach \s in {0.35,1.175,2.0,2.825,3.65}{ \foreach \z in {0.35,1.3,2.25}{ \fill[perno] (\s,\z) circle (1.6pt); } }
    \node[font=\small, anchor=north] at (\W/2,-0.1) {\textbf{B}: grilla abierta};
  \end{scope}
  % C: tresbolillo
  \begin{scope}[shift={(9.2,0)}]
    \fill[acero!15] (0,0) rectangle (\W,\Hp); \draw[thick] (0,0) rectangle (\W,\Hp);
    \foreach \z [count=\k] in {0.3,0.8,1.3,1.8,2.3}{
      \pgfmathsetmacro{\off}{mod(\k,2)==0 ? 0.4 : 0}
      \foreach \s in {0.3,1.1,...,3.5}{ \fill[perno] ({\s+\off},\z) circle (1.6pt); } }
    \node[font=\small, anchor=north] at (\W/2,-0.1) {\textbf{C}: tresbolillo};
  \end{scope}
  % D: densificada en esquinas y extremos
  \begin{scope}[shift={(13.8,0)}]
    \fill[acero!15] (0,0) rectangle (\W,\Hp); \draw[thick] (0,0) rectangle (\W,\Hp);
    \foreach \s in {0.25,0.6,1.05,1.65,2.35,2.95,3.4,3.75}{
      \foreach \z in {0.25,0.6,1.3,2.0,2.35}{ \fill[perno] (\s,\z) circle (1.6pt); } }
    \node[font=\small, anchor=north] at (\W/2,-0.1) {\textbf{D}: densa en bordes};
  \end{scope}
\end{tikzpicture}
\caption{Las cuatro distribuciones de pernos a probar, en una plancha vista desde el hormig\'on
(horizontal: a lo largo de la pared, de esquina a esquina; vertical: eje del t\'unel). Cada una se
corre con varios espesores $t$.}
\label{fig:distrib}
\end{figure}

%================================================================
\section{Metodolog\'ia}
%================================================================

Se usa un c\'odigo propio de elementos finitos 3D en Python que extiende el solver 2D de las tareas
anteriores.
\begin{enumerate}
  \item \textbf{Elementos:} hexaedros H8, H20 y H27, cuyas caras son los Q4, Q8 y Q9 del c\'odigo 2D. En
        el blindaje se usa H20 con un elemento en el espesor (H8 sufre \emph{shear locking}).
  \item \textbf{Malla param\'etrica:} pasa por las esquinas, uniones de planchas y cada perno, que cae
        en un nodo; cambiar $t$, $p_s$, $p_z$ o $e$ la regenera sola.
  \item \textbf{Condiciones de borde:} los pernos son las condiciones de borde del modelo: en cada
        nodo de perno se impone $\mathbf u = 0$.
  \item \textbf{Tensi\'on junto al perno:} el apoyo puntual crea una singularidad, as\'i que se eval\'ua
        $\sigma_c$ en el centro del pa\~no y $\sigma_p$ a distancia fija del perno (borde de soldadura).
  \item \textbf{Barrido:} se corren las 4 distribuciones con cada espesor, se tabulan tensiones,
        reacciones y flechas, y se combinan la mejor distribuci\'on y el menor espesor que cumple.
\end{enumerate}

%================================================================
\section{Resultados esperados}
%================================================================

Curvas de tensi\'on, reacci\'on en el perno m\'as cargado y flecha en funci\'on de $t$ para cada
una de las cuatro distribuciones; una tabla de qu\'e combinaciones cumplen; la comparaci\'on del
\'optimo con el dise\~no original (ahorro de acero y pernos, margen frente a los admisibles); y los mapas de Von Mises y la
deformada del dise\~no \'optimo, identificando las zonas cr\'iticas (esquinas, extremos y pernos de
borde).

\end{document}
